% Capítulo 5: Repaso

\section*{CAPÍTULO 5 \quad Repaso}

\subsection*{Temas Básicos}

\noindent \textit{Definir, establecer, o analizar lo siguiente.}
\begin{enumerate}
    \begin{minipage}{0.45\textwidth}
    \item Área entre curvas
    \item Volumen de un sólido general
    \item Volumen de un sólido de revolución: método de los discos (o arandelas)
    \item Volumen de un sólido de revolución: método de las envolventes cilíndricas.
    \end{minipage}%
    \begin{minipage}{0.45\textwidth}
    \item Trabajo
    \item Valor medio de una función
    \item Teorema del Valor Medio para Integrales
    \end{minipage}
\end{enumerate}

\subsection*{Ejercicios}

\noindent \textit{Calcule el área de la región limitada por las curvas dadas en cada uno de los Ejercicios 1 a 8.}

\begin{enumerate}
    \begin{minipage}{0.45\textwidth}
    \subimport{ejercicio_01/}{ejercicio.tex}
    \subimport{ejercicio_03/}{ejercicio.tex}
    \subimport{ejercicio_05/}{ejercicio.tex}
    \subimport{ejercicio_07/}{ejercicio.tex}
    \end{minipage}%
    \begin{minipage}{0.45\textwidth}
    \subimport{ejercicio_02/}{ejercicio.tex}
    \subimport{ejercicio_04/}{ejercicio.tex}
    \subimport{ejercicio_06/}{ejercicio.tex}
    \subimport{ejercicio_08/}{ejercicio.tex}
    \end{minipage}
\end{enumerate}

\subsection*{Ejercicios 9--13: Rotaciones alrededor de los ejes cartesianos.}

\noindent \textit{Calcule el volumen del sólido que se obtiene al hacer girar en torno al eje indicado la región limitada por las curvas dadas en cada uno de los Ejercicios 9 a 13.}

\begin{enumerate}
    \setcounter{enumi}{8}
    \subimport{ejercicio_09/}{ejercicio.tex}
    \subimport{ejercicio_10/}{ejercicio.tex}
    \subimport{ejercicio_11/}{ejercicio.tex}
    \subimport{ejercicio_12/}{ejercicio.tex}
    \subimport{ejercicio_13/}{ejercicio.tex}
\end{enumerate}

\subsection*{Ejercicios 14--16: Planteamiento de integrales para volúmenes.}

\noindent \textit{Establezca, pero no evalúe, una integral para el volumen del sólido que se obtiene al hacer girar en torno al eje indicado la región limitada por las curvas dadas en cada uno de los Ejercicios 14 a 16.}

\begin{enumerate}
    \setcounter{enumi}{13}
    \subimport{ejercicio_14/}{ejercicio.tex}
    \subimport{ejercicio_15/}{ejercicio.tex}
    \subimport{ejercicio_16/}{ejercicio.tex}
\end{enumerate}

\subsection*{Ejercicios 17--18: Áreas y volúmenes combinados de regiones complejas.}

\begin{enumerate}
    \setcounter{enumi}{16}
    \subimport{ejercicio_17/}{ejercicio.tex}
    \subimport{ejercicio_18/}{ejercicio.tex}
\end{enumerate}

\subsection*{Ejercicios 19--22: Identificación y descripción de sólidos definidos por integrales dadas.}

\noindent \textit{Las integrales dadas en los Ejercicios 19 a 22 representan volúmenes de sólidos. Describa los sólidos.}

\begin{enumerate}
    \setcounter{enumi}{18}
    \begin{minipage}{0.45\textwidth}
    \subimport{ejercicio_19/}{ejercicio.tex}
    \subimport{ejercicio_21/}{ejercicio.tex}
    \end{minipage}%
    \begin{minipage}{0.45\textwidth}
    \subimport{ejercicio_20/}{ejercicio.tex}
    \subimport{ejercicio_22/}{ejercicio.tex}
    \end{minipage}
\end{enumerate}

\subsection*{Ejercicios 23--32: Sólidos transversales, monumentos, cono modular, resortes, vaciado de paraboloides de revolución y límites.}

\begin{enumerate}
    \setcounter{enumi}{22}
    \subimport{ejercicio_23/}{ejercicio.tex}
    \subimport{ejercicio_24/}{ejercicio.tex}
    \subimport{ejercicio_25/}{ejercicio.tex}
    \subimport{ejercicio_26/}{ejercicio.tex}
    \subimport{ejercicio_27/}{ejercicio.tex}
    \subimport{ejercicio_28/}{ejercicio.tex}
    \subimport{ejercicio_29/}{ejercicio.tex}
    \subimport{ejercicio_30/}{ejercicio.tex}
    \subimport{ejercicio_31/}{ejercicio.tex}
    \subimport{ejercicio_32/}{ejercicio.tex}
\end{enumerate}
